\section{Resultados Parciais e Esperados}
\label{sec:resultados}

Esta seção descreve os avanços já realizados na pesquisa, validados por meio de publicação científica, e detalha os resultados finais que se espera alcançar com a aplicação completa da metodologia proposta.

\subsection{Resultados Parciais}
Como prova de conceito da etapa de modelagem territorial e cálculo de indicadores, foi desenvolvido o estudo intitulado \textit{"Fatores Socioeconômicos e Desigualdade na Mobilidade de Curta Distância aos Serviços Urbanos de Curitiba"} \cite{sbbd_estendido}, apresentado no Apêndice~\ref{apendice:fatores_socioeconomicos}. O trabalho, focado em deslocamentos a pé, estruturou a rede viária e os Pontos de Interesse das categorias de alimentação, comércio, cultura, educação, saúde e transporte, e calculou métricas de acessibilidade como o PTh, PTc e F15 para o cenário típico.

Por meio de análise de regressão e clusterização espacial, o estudo cruzou os tempos de acesso com variáveis socioeconômicas. Os resultados mostraram que variáveis como escolaridade, empregabilidade e acesso à infraestrutura básica estão diretamente associadas a melhores tempos de acesso aos serviços urbanos em Curitiba. Áreas com maior concentração de pessoas com mais anos de estudo ou maior frequência ao ensino médio tendem a apresentar tempos de deslocamento mais baixos e acesso facilitado, enquanto regiões marcadas por maior vulnerabilidade social, como presença de jovens fora da escola e do trabalho, precariedade de água e esgoto, e moradias em condições inadequadas, registram deslocamentos mais longos. 

A análise espacial confirmou que os maiores tempos de acesso se concentram nos setores periféricos da cidade, ampliando desigualdades já presentes nos indicadores socioeconômicos. Além disso, o estudo também identificou casos pontuais de territórios (\textit{outliers}) com oferta de serviços acima da média do entorno imediato, assim como áreas mais próximas as regiões centrais em situação de exclusão devido à carência de infraestrutura urbana.

Embora o estudo tenha consolidado a coleta de dados e o cálculo base da acessibilidade, ele apresentou limitações que fundamentam o avanço metodológico desta dissertação. O modelo utilizou uma topologia estática, focando exclusivamente no cenário típico a pé, sem considerar barreiras físicas temporárias, perturbações ambientais ou a qualidade e capacidade de atendimento dos serviços.

\subsection{Resultados Esperados}

Com a aplicação integral da metodologia proposta, espera-se demonstrar de forma quantitativa como perturbações climáticas e ambientais podem degradar significativamente a acessibilidade urbana em comparação ao cenário típico. O trabalho deverá validar o mecanismo de penalização do grafo, mostrando na prática como a introdução de camadas ambientais altera o custo do deslocamento e modifica a topologia funcional da cidade. Os resultados esperados incluem a produção de mapas detalhados, capazes de identificar regiões que, embora tradicionalmente bem servidas de acessibilidade, se tornam isoladas ou apresentam forte degradação sob condições ambientais adversas.

Além disso, a análise irá quantificar os efeitos dessas perturbações sobre diferentes grupos sociais, por meio da aplicação de métricas como o Índice G e o cruzamento com indicadores socioeconômicos. Assim, será possível evidenciar se as perdas de acesso afetam a população de forma homogênea ou se recaem de maneira mais intensa sobre territórios e grupos em situação de vulnerabilidade. Espera-se, portanto, fornecer um panorama analítico robusto que possa subsidiar o planejamento urbano e orientar políticas públicas mais resilientes, considerando não apenas a oferta de oportunidades, mas também as barreiras dinâmicas impostas pelo clima e pelo ambiente ao acesso dos cidadãos.